\section{Intrinsic observations and the rigidity theorem}
\label{sec:introduction}
Near the first possible occurrence of a collision itinerary, its
probability integrates both the action of the broken ray and the physical
flux through a collision section. We study the geometric information in
this law. The central observation is intrinsic: endpoint positions are
recorded by distance along the reflecting boundary, not by a supplied
Cartesian contact chart. A single reversal of all the arclength signs
is immaterial. The resulting inverse recovers the leading geometry as
well as the higher asymmetric jets.

There are two distinct steps from a local law to table rigidity.
First, a finite collection of exact moment germs and limits determines
a participating pair of analytic obstacles. These germs contain
infinitely many coefficients. Second, overlapping pair reconstructions
must be synchronized on a connected channel network. Intrinsic contact
registration performs this step for arbitrary analytic obstacles.
For asymmetric obstacles, however, equality of the entire reconstructed
images supplies the synchronization itself: independent edge reversals
and unknown contact offsets need not be coordinated in advance.
Lattice-cycle marking then determines the periods. Quantitative analytic
and asymmetry priors also permit finite noisy-window reconstruction of
the whole table, with the onset estimated as part of the experiment.
No claim is made that an unvisited obstacle is determined by one channel.

\subsection{Geometric and observational conventions}
Consider finitely many labelled compact strictly convex obstacles modulo
a full-rank lattice in $\R^2$. Their boundaries are smooth with positive
curvature, their closures are disjoint, and the free cell area is $A>0$.
Initial conditions have normalized phase volume
$\dd q\dd\vartheta/(2\pi A)$ and speed one. A selected channel is an
isolated unobstructed common normal segment between two designated
obstacle copies. Its endpoints are labelled $0,1$. Fixed sufficiently
small facing patches isolate this channel; clearance and positive
separation guarantee a common small time collar.

The two endpoint origins are part of the channel marking. Their ambient
coordinates, the normal direction, the gap, the curvatures, and the area
are not supplied numerically. Choose either orientation of the plane
solely to describe the observation. Let $n$ point from contact $0$ to
contact $1$ and let $t$ be its positive quarter-turn. At both endpoints
use signed arclength increasing in tangent direction $t$. Equivalently,
one may use counterclockwise arclength around each obstacle and reverse
its sign at contact $1$. This prescription involves no numerical
Cartesian chart. Reversing the plane orientation reverses both signs.
The same choice is used for both itineraries and all windows.

For a time window $T$, let $E_b(T)$ be the event of exactly the three
successive impacts $b,1-b,b$ in the selected patches. On success record
the first and last signed boundary arclengths $(s_-,s_+)$ from contact
$b$; failures have a separate symbol. Define the subprobability measure
\begin{equation}\label{eq:intrinsicmeasure}
 \mu_{b,T}(B)=\Prob\{E_b(T),\ (s_-,s_+)\in B\},\qquad B\subset\R^2.
\end{equation}
The intrinsic observation is the orbit of the entire family
$(\mu_{0,T},\mu_{1,T})_T$ under
$\iota(s_-,s_+)=(-s_-,-s_+)$. It is \emph{not} the average of a
measure with its reflection, nor a separate sign choice at each time.
In particular, the observation retains odd information while supplying
no preferred orientation. It uses a boundary-arclength sensor and is
strictly richer than a collision count.

Write $T_0$ for the local onset of positive mass and $d=T-T_0$.
From any one representative of the observation extract
\begin{align}
 P_b(d)&=\mu_{b,T_0+d}(\R^2),&
 M_b(d)&=\int(s_-+s_+)\dd\mu_{b,T_0+d},\label{eq:arcdatum}\\
 D_b(d)&=\int(s_--s_+)^2\dd\mu_{b,T_0+d},&
 \beta_b&=\lim_{d\downarrow0}\frac{D_b(d)}{dP_b(d)}.\label{eq:betadef}
\end{align}
Thus the proof below uses only $T_0$, two mass germs, two first-moment
germs, and two numbers $\beta_b$, not the whole marked distribution.
The moments are unconditional; no observation or preparation is discarded
on failure. The quotient of these quantities by their single common sign
is a canonical finite collection of exact germs and limits extracted
from \eqref{eq:intrinsicmeasure}. It is not a finite list of real
numbers or a finite-sample observation.

\begin{theorem}[Intrinsic local and global reconstruction]
\label{thm:intrinsic}
The following assertions hold for the observation above.

\smallskip\noindent
\textup{(i)} Its onset and quadratic moments recover the leading geometry:
\begin{equation}\label{eq:calibrationmain}
 g=\frac{T_0}{2},\qquad
 \kappa_b=\frac{2}{3\beta_b}-\frac1g,\qquad
 A=\frac1{4\alpha\sqrt{c_0c_1(c_0c_1-1)}},
\end{equation}
where $c_b=1+g\kappa_b$ and
$\alpha=\lim_{d\downarrow0}P_b(d)/d^2$ is independent of $b$.

\smallskip\noindent
\textup{(ii)} At every fixed order $K\ge3$, these leading data and the
coefficients of $P_b/(\alpha d^2)$ and $M_b/(\alpha d^2)$ through the
orders specified in Theorem~\ref{thm:main} determine the contact $K$-jets
up to simultaneous transverse reflection. On either coherent orientation
cover the inverse is analytic and locally bi-Lipschitz at finite order.
Its highest-jet blocks are exactly \eqref{eq:evenblock} and
\eqref{eq:oddblock}; arbitrary lower odd jets and equal curvatures are
allowed. For connected closed real-analytic obstacle boundaries, the
observation determines the entire participating pair, in its relative
position, up to one Euclidean isometry.

\smallskip\noindent
\textup{(iii)} Suppose every obstacle is visited by a connected finite
registered normal-channel network as defined in
Section~\ref{sec:global}. If the marked lattice displacements of its
cycles span $\R^2$, its intrinsic observations and registration determine
the whole labelled periodic table and its marked lattice up to one
Euclidean isometry. No obstacle image or numerical contact frame is
supplied.

\smallskip\noindent
\textup{(iv)} On a connected full-rank network whose analytic obstacles
have trivial Euclidean isometry groups, the contact registration and
relative edge signs in (iii) need not be supplied. Independent local
reversal orbits determine the same labelled table and marked lattice.
If all obstacles are noncircular, these unregistered data leave only
finitely many possible assemblies. The precise bound and a nonempty
open physical asymmetric class are given in Section~\ref{sec:unregistered}.

\smallskip\noindent
\textup{(v)} On the uniformly analytic classes separated from symmetry
specified in Section~\ref{sec:stable}, finitely many noisy moment
windows give a convergent regularized reconstruction of the entire table.
For each fixed $K\ge3$ its error is bounded by
\[
 C\bigl(A_K\epsilon^{1/(\lceil K/2\rceil+2)}+Bq^K\bigr)^\vartheta,
       \qquad 0<q,\vartheta<1,
\]
when the scalar mean errors are at most sufficiently small $\epsilon$.
The onset and leading geometry are estimated, not supplied exactly.
The constants $A_K$ need not be uniform in $K$.
\end{theorem}

Parts (i)--(iv) concern exact germs. The finite-order stability in
(ii) is not stability of analytic continuation, is not uniform in $K$,
and does not turn an onset limit into a finite-sample measurement.
Part (iii) includes explicit intrinsic registration and lattice labels;
these are additional data, not consequences of a single channel law.
At reflection-fixed jets the unframed space is a finite quotient rather
than a smooth coordinate manifold. The asserted analytic coordinates
are on a coherent orientation cover.

\subsection{Cartesian coordinates used in the proof}
After the gap has been identified, put the two contacts at $(0,0)$ and
$(g,0)$. The facing boundaries have inward height functions
\begin{equation}\label{eq:graphs}
 \psi_b(y)=\frac{\kappa_by^2}{2}
       +\sum_{r\ge3}\frac{q_{b,r}}{r!}y^r,\qquad \kappa_b>0.
\end{equation}
The actual positions are $(-\psi_0(y),y)$ and
$(g+\psi_1(y),y)$. For smooth boundaries the series denotes jets only.
The corresponding signed arclength is
\begin{equation}\label{eq:arcmap}
 \mathfrak a_b(y)=\int_0^y\sqrt{1+\psi_b'(x)^2}\dd x
       =y+\frac{\kappa_b^2y^3}{6}+O(y^4).
\end{equation}
Hence $M_b$ is the expectation of $\mathfrak a_b(u)+\mathfrak a_b(v)$ on success.
The auxiliary Cartesian moment and normalized coefficients are
\begin{gather}
 R_b(d)=\E[(u+v)\one_{E_b(2g+d)}],\qquad
 E_{2,b}(d):=E_b(2g+d),\label{eq:rawdata}\\
 z=c_0c_1-1,\quad \gamma=\arcosh\sqrt{c_0c_1},\quad
 \alpha=\frac1{2A\sinh(2\gamma)},\label{eq:leading}\\
 G_b(d)=\frac{P_b(d)}{\alpha d^2}=1+\sum_{k\ge1}\xi_{b,k}d^k,
 \quad J_b(d)=\frac{R_b(d)}{\alpha d^2}=\sum_{k\ge1}\eta_{b,k}d^k.
 \label{eq:coefficients}
\end{gather}
The Cartesian moment is used to compute universal blocks; it is not
silently substituted for the intrinsic observable. Lemma~\ref{lem:filter}
proves the precise transfer to arclength.

\begin{theorem}[Cartesian signed-moment component]\label{thm:main}
At supplied labelled gap, separate curvatures and common transverse
orientation the following assertions hold.
\textup{(i)} The common leading probability coefficient is $\alpha$ and
determines $A$. \textup{(ii)} For
$n_e=\lfloor K/2\rfloor-1$, $n_o=\lfloor(K-1)/2\rfloor$,
the coefficient vectors $(\xi_{b,k})_{1\le k\le n_e}$ and
$(\eta_{b,k})_{1\le k\le n_o}$, $b=0,1$, are analytic triangular
coordinates for $(q_{b,r})_{3\le r\le K}$, locally bi-Lipschitz at
fixed order on compact positive leading-geometry boxes.
\textup{(iii)} For analytic contacts the four raw germs $P_b,R_b$
determine the participating analytic boundary images in these frames.
\textup{(iv)} Analytic periodic physical families realize all the displayed
jet directions, at fixed area or with area free. Respectively $2K-4$ or
$2K-3$ positive-window scalar means are local bi-Lipschitz coordinates
on these supplied families. At fixed area their selected leading
count and endpoint-metric hierarchy is fixed at every flight number.
\end{theorem}

We retain the full Cartesian proofs, physical realization and charged
binary experiment in Sections~\ref{sec:local}, \ref{sec:inverse} and
\ref{sec:physical}. The standalone filtration below supplies the
arbitrary-lower-jet justification, and Lemma~\ref{lem:analyticidentity}
supplies the full continuation statement used in part (iii).

\subsection{Relation to inverse billiard problems}
B\'alint, De Simoi, Kaloshin and Leguil~\cite{BDKL} recover period-two
curvatures and periodic Lyapunov exponents from marked lengths of open
dispersing billiards. De Simoi, Kaloshin and Leguil~\cite{DSKL} obtain
analytic table determination under symmetry and genericity assumptions.
Osterman~\cite{Osterman} relates marked lengths to homoclinic analytic
conjugacy and determines a third scatterer from two known scatterers and
the marked spectrum. Finamore and Leguil~\cite{FL} prove isometric
rigidity from an enriched marked length spectrum for finite-horizon
Sinai billiards. These are global dynamical observations; none is
identified here with a near-onset arclength law.

The distinction of mechanism also matters. Koval~\cite{Koval} studies
formal billiard reconstruction and Gevrey divergence with odd derivatives
as parameters. Our highest-jet calculation is a physical residual-flux
integral, not an assertion that every formal inverse series converges.
De Simoi, Kaloshin and Wei~\cite{DSKW} prove deformational dynamical
spectral rigidity near a circle in a symmetric convex class. Their
periodic-length observation and focusing geometry differ from those
used here. We make no assertion of exhaustive priority from this
comparison.

The present result removes supplied leading geometry by a quadratic
moment calibration and replaces Cartesian marks by intrinsic metric
marks. It joins local determinations either by finite registration or,
for asymmetric obstacles, by unique matching of the recovered analytic
images. The latter mechanism does not insert cross-channel gluing
information into the observation. A further quantitative argument controls
finite coefficient recovery, analytic continuation, image matching, and
lattice inversion separately. Conditional continuation is necessarily
sensitive to its analytic prior; our use of the two-constants principle
is consistent with the continuation estimates discussed by
Trefethen~\cite{Trefethen}, and claims no optimal full-inverse rate.

The observation has an intrinsic geometric origin: it is the pushforward
of normalized Liouville measure by the selected itinerary and its
boundary-distance marks. A Euclidean isometry preserves that measure,
the itinerary and boundary distance, acting only by the stated reversal
when orientation changes. The selection and local contact origins are
still marks; this invariant is richer than counts and is not identified
with a marked length spectrum. The unregistered theorem reduces its
inter-experiment information, rather than claiming equivalence to a
conventional global spectrum. It does
not claim count-only asymmetric rigidity. Section~\ref{sec:countfibers}
describes the finite and formal count-data fibers, while distinguishing
these from the exact analytic count-germ question. The smooth relative
boundary law and its function-valued Volterra and Abel theory retain their
own information category in Section~\ref{sec:relative} and Supplement S.
